
\(\underline{E}_1,\ldots,\underline{E}_r\) de \(\underline{O}\)-modules localement libres de rang \(\underline{n}_1,\ldots,\underline{n}_r\)) par recollement du \(\underline{O}\)-module libre \(T(\underline{O}^{\underline{n}_1},\ldots,\underline{O}^{\underline{n}_r})\) sur des cartes locales trivialisantes moyennant la représentation \(\alpha\). De façon précise, à la suite \(\underline{E}_1,\ldots,\underline{E}_r\) correspond un \(G\)-torseur à droite \(P=P_{\underline{E}_1,\ldots,\underline{E}_r}\) tel que l'on a un isomorphisme canonique
\begin{equation}\tag{6.1.3}
 \underline{E}_i\xrightarrow{\sim}P\overset{G}{\wedge}\underline{O}^{\underline{n}_i}\qquad(i=1,\ldots,r),
\end{equation}
où \(\overset{G}{\wedge}\) désigne le produit contracté, et \(G\) opère sur \(\underline{O}^{\underline{n}_i}\) via la \(i\)-ième projection \(G\to\Gl(\underline{n}_i)\) ; alors toute représentation
\[
 \alpha:G\longrightarrow\underline{\Aut}_{\underline{O}}(M)
\]
fait opérer \(G\) sur \(\underline{O}^{\underline{m}}\) par des \(\underline{O}\)-automorphismes, et définit ainsi un multifoncteur \(T=T_\alpha\) tel que
\begin{equation}\tag{6.1.4}
 T_\alpha(\underline{E}_1,\ldots,\underline{E}_r)=P_{\underline{E}_1,\ldots,\underline{E}_r}\overset{G}{\wedge} M.
\end{equation}
De plus, pour toute représentation composée \(\alpha=\alpha_1\circ\alpha_2\) on trouve
\begin{equation}\tag{6.1.5}
 T_\alpha=T_{\alpha_1}\circ T_{\alpha_2}.
\end{equation}

\subsection*{6.2.}
Pour tout \(\underline{O}\)-module localement libre \(\underline{E}\) de rang \(\underline{n}\) nous désignerons par \(\det(\underline{E})\) le \(\underline{O}\)-module inversible \(\bigwedge\limits^{\underline{n}} \underline{E}\). Gardant les notations de 6.1, le foncteur \(\det(T(\underline{E}_1,\ldots,\underline{E}_r))\) est déterminé, à isomorphisme canonique près, par la représentation correspondante
\begin{equation}\tag{6.2.1}
 \det\circ\alpha_T:G\longrightarrow{G}_m=\Gl(1).
\end{equation}

Supposons maintenant que \(X\) soit le topos annelé des faisceaux sur le site annelé \((\Sch/S)\) muni de la topologie étale. Alors on peut calculer \(\det(T(\underline{E}_1,\ldots,\underline{E}_r))\), ou ce qui revient au même, la représentation \(\det\circ\alpha_T\), à l'aide de la

